\scanpage{4}

% L1–7, continuation on the left
existe un drapeau majorant $x$ et $y$ ? \unsure{En fait|Il faut|fois} \unsure{donc|dans} \inserted{l'ens ordonné} $F_*$ se récupère via l'ens ordonné des drapeaux (qui a la propriété, entre autres, que tout élément est le sup des éléments minimaux qu'il majore, \inserted{\deleted{\unsure{à is. près avec son ordre}} et des \unsure{paires|l'ens.} des drapeaux \unsure{maj. par} un drapeau} \uuline{cf?).} Drapeaux maximaux\ldots{} Soit $\mathcal A$ l'ens des drapeaux maximaux,

% L1–4, extreme left
% slanted
\margin{au passage à l'ordre inverse près \ldots}

% L8–9, F(x)
\deleted{\unsure{$F_\alpha$}} Pour tout $x\in F_{\deleted{*}\inserted{\alpha}}$, soit $F(x)$ l'ens des $y<x$. \inserted{Il satisfait les m\textsuperscript{es} conditions que $F_*$, avec $n$ remplacé par $\alpha-1$ si $\alpha$\unsure{$\geq 1$|$\neq0$}.} \deleted{On posera l'axiome}

% L10–11, first numbered condition
\uline{Axiome 2} Soit \inserted{$2\leq\alpha\leq n$}, $x_\alpha\in F_\alpha$, $x_{\alpha-2}\in F(x_{\unsure{\alpha|{}}})_{\alpha-2}$. Alors \deleted{\unsure{Il}} le cardinal de l'ens des $y_{\alpha-1}\in F(x_{\unsure{1|\alpha}})_{\alpha-1}$ incidents à $x_{\alpha-2}$ est 2.

% L10–11, left bracket
% slanted
\margin{\unsure{admettre|à mettre|condition} seulement si $n\geq2$}

% L12–13, second condition
\uline{Axiome 3} Soit $x_1\in F_1$, alors le cardinal de l'ens des $y_0\in\deleted{\unsure{F_0,\ x}}\ F(x)_0$ \deleted{incidents à $x$} est 2.

% L12–13, second left bracket
% slanted
\margin{à \unsure{mettre} si $n\geq1$}

% L14–18, maximal flags
\deleted{\unsure{Définissons|Définir|Définition}} Soit alors $\mathcal A$ l'ens des drapeaux maximaux. \deleted{\unsure{de G}} On va introduire sur $\mathcal A$ des opérations $\sigma_\alpha$ $(0\leq\alpha\leq n\deleted{\unsure{\}}}-1)$. Soit $d=(x_0,x_1,\ldots,x_\nu)$ $(\nu\leq n)$ un drapeau maximal, avec $x_i\in F_i$ $(0\leq i\leq\nu)$.

% L19–21, first operator
\[
\begin{aligned}
\sigma_0(d)&=(x'_0,x_1,\ldots,x_\nu)\qquad\deleted{\text{où }x'_0\text{ est}}\ \text{caractérisé par la condition}\\
&\hspace{7em}x'_0\ne x_0\ \big[\text{\uline{et}}\quad(x'_0,x_1,\ldots x_\nu)\ \text{un drapeau}\\
&\hspace{7em}\text{i.e. }x'_0\in F(x_1)_0\big]\qquad\text{cf Axiome 2}
\end{aligned}
\]

% L22–23, remaining operators
\[
\begin{aligned}
1\leq\alpha\leq n-1\quad \sigma_\alpha(d)&=(x_0,\ldots x_{\alpha-1},x'_\alpha,x_{\alpha+1},\ldots)\quad\text{caractérisé par }x'_\alpha\ne x_\alpha\\
&\hspace{8em}\text{\uline{et}}\ \text{que }\sigma_\alpha(d)\ \text{est un drapeau.})
\end{aligned}
\]

% L24–26, optional axiom
Considérons aussi l'axiome (facultatif)

\uline{Axiome $3'$} $\forall x_{n-1}\in F_{n-1}$, le cardinal de l'ens des \deleted{drapeaux} $y\in F_n$ incidents à $x_{n-1}$ est 2.

% L25–27, outer left bracket
% slanted
\margin{pour $n\geq1$ fixé}

% L27–28, last operator
On pose alors
\[
\sigma_n(d)=(x_0,\ldots,x_{n-1},x'_n),\quad \deleted{\text{où }x'_n\text{ est le}}\ \inserted{\text{caractérisé par }x'_n\ne x_n\quad\text{\unsure{(et\ldots)|(etc.\ldots)}}}
\]

% L29–30, foot of left column
\uline{NB} Les $F(x)$ (où $x\in F_\alpha$) satisfont les axiomes $1,2,3$, \uline{et $4$} avec $n$ remplacé par $\alpha-1$\ldots

% R1–3, head of right column
On va voir maintenant que (moyennant les axiomes 1, $1'$, 2, \unsure{3, $3'$|$2'$, 3} \inserted{et un axiome 4 plus bas}) l'ens à opérateurs $\mathcal A$ permet de \uline{reconstituer} $F_*$ (donc $(F_\alpha),(I_\alpha)$\ldots).

% R4–6, commutations
\deleted{On \unsure{vérifie|a évidemment}} \deleted{Les} Mais donnons d'abord les relations de commutation évidentes
\[
 \sigma_\alpha\sigma_\beta=\sigma_\beta\sigma_\alpha\qquad\text{si }|\beta-\alpha|\ne1\quad\text{i.e. }\alpha,\beta\ \text{non consécutifs.}
\]

% R7–11, group and subgroup
Soit $G_n$ le groupe engendré par les \deleted{\unsure{rela\textsuperscript{ts}|relations}} gén. $\sigma_\alpha$ $(0\leq\alpha\leq n)$ avec les relations précédentes. $\mathcal A$ est un $G_n$-ensemble. Les $\sigma_\alpha$ y opèrent sans pt fixe. Soit $G_n^{(\alpha)}$ \inserted{$(0\leq\alpha\leq n)$} le sous-groupe de $G_n$ engendré par les $\sigma_\beta$ avec $\beta\ne\alpha$.

% R12–13, orbit map
On a évidemment
\[
 \mathcal A/G_n^{(\alpha)}\longrightarrow F_\alpha\qquad\unsure{(\text{si }\alpha\leq n-1)|(0\leq\alpha\leq n-1)}
\]

% R14–20, fourth axiom with justification
\deleted{\unsure{Proposition}} \inserted{\uline{Axiome 4}} Cette application est bijective. (La surjectivité est claire \inserted{par l'axiome 1}. Pour la bijectivité, il suffit de voir que si $d=(x_0,\ldots x_n)$, $d'=(\unsure{x_0|x'_0},\ldots x'_n)$ sont des él\textsuperscript{ts} avec \deleted{$d,d'$ \illegible{1}}, $x_\alpha=x'_\alpha$, alors \deleted{\unsure{op}} on peut passer de $d$ à $d'$ en appliquant successiv\textsuperscript{mt} des transformations $\sigma_\beta$ avec $\beta\ne\alpha$.

% R14–15, in the gutter
% slanted
\margin{le \unsure{rendre} naturel}

% R21–24, a general type
Plus généralement, soit $\delta\subset[0,n]$, \inserted{$\delta\ne\varnothing$} et considérons le sous-groupe $G_n^{(\delta)}$ de $G_n$, engendré par les $\sigma_i$ avec $i\in\complement_{[0,n]}\delta$, Alors on a

% R25, lower right
\[
 \mathcal A/G_n^{(\unsure{\alpha|\delta})}\unsure{-|\to}\mathcal A_\delta\quad\text{(drapeaux de type }\delta\text{)}
\]

% R26–28, final question
et on voudra que ce soit bijectif (est-ce une conséquence de l'axiome 4 et des précédents?).

% bottom left
% archive sheet number: 3
